\documentclass[conference]{IEEEtran}
\usepackage{amsmath}
\usepackage{cite}
\usepackage{hyperref}

\begin{document}

\title{Data-Driven Priors and Physics-Based Corrections: A Hybrid Flow
Matching Approach to Safe Robot Motion Planning}

\author{
  \IEEEauthorblockN{Zihao Sun}
  \IEEEauthorblockA{606466419}
  \and
  \IEEEauthorblockN{Charles Beck}
  \IEEEauthorblockA{406045359}
}

\maketitle

\begin{abstract}
  UCLA ECE M237 Spring 2026 Course Project Milestone Reports. Week
  2---Literature Review.
\end{abstract}

% -------------------------------------------------------
\section{Main Idea}

We investigate how to merge an offline-trained flow matching trajectory
prior (learned from demonstration data) with an online model-based score
correction (computed analytically from the known robot dynamics, cost
function, and safety constraints), and study whether the resulting
hybrid produces trajectories that are simultaneously diverse, physically
feasible, and safe.

% -------------------------------------------------------
\section{Literature Overview}

Our project sits at the intersection of two primary bodies of
literature, with a secondary connection to a third.
The first is \emph{model-based score computation}: Model-Based
Diffusion (MBD)~\cite{mbd} derives the score analytically from
known dynamics and cost, guaranteeing feasibility without any learned
component; its authors note data ``can be naturally integrated'' but
leave this unformalized---the gap we address.
The second is \emph{generative trajectory planning}: diffusion-based
planners such as Diffuser~\cite{janner2022} and flow-matching
methods such as FlowMP~\cite{flowmp} learn expressive
multimodal priors from demonstrations but provide no feasibility
guarantees out of distribution; kernel-based Behavioral Score
Diffusion~\cite{bsd} offers a training-free data-driven variant.
As a secondary body of literature, we review work on \emph{safe
generative control} via CBFs~\cite{safeflow, s2diff},
which we plan to incorporate as an independent score term for collision
avoidance.  Further connections include classifier-free guidance as
policy improvement~\cite{cfgrl} and optimal-control-guided flow
matching~\cite{ocflow}.

% -------------------------------------------------------
\section{Highlighted Literature Review}

\subsection{Summary}

Pan et al.~\cite{mbd} propose \emph{Model-Based Diffusion}
(MBD), a framework that replaces the learned score of standard
diffusion planners with one computed analytically from the robot's
dynamics model $x_{k+1}=f(x_k,u_k)$ and cost $J(\tau)$ via Tweedie's
formula and self-normalized importance sampling.
Concretely, MBD defines a Boltzmann trajectory distribution
$p^{\mathrm{mod}}(\tau)\propto\exp(-J(\tau)/\lambda)$, noises it via a
VP-SDE, and approximates the posterior mean by cost-reweighted samples.
On Maze2D, Block Push, and Franka manipulation benchmarks, MBD matches
or outperforms Diffuser~\cite{janner2022} while requiring zero
demonstration data.

\subsection{Key Idea \& Relevance}

The central insight of MBD is that for systems with known dynamics the
score $\nabla \log p^{\mathrm{model}}_t$ can be computed exactly by
differentiating through the dynamics and cost at each sampling step,
yielding hard feasibility that learned methods cannot guarantee.
However, MBD's authors note this score ``can be naturally integrated''
with data-driven priors yet provide no formalization.
Our project takes this as its central question, proposing to fuse the
model-based score with a flow matching prior via a product-of-experts:
\begin{equation}
  \nabla \log p^{\mathrm{target}}_t \;\approx\;
  \nabla \log p^{\mathrm{data}}_t
  + \lambda(x, t)\,\nabla \log p^{\mathrm{model}}_t,
\end{equation}
where the state- and time-dependent weight $\lambda(x,t)$ is designed
to trust the data prior in regions well-covered by the training
distribution, and to rely on the physics model in out-of-distribution
states.
This adaptive weighting mechanism---and the empirical study of the
tradeoff it governs---is the primary novel contribution of our work.
We further conjecture that in the linear-quadratic special case the
model-based correction analytically reduces to the LQR gain
$-\lambda^{-1}P_k x_k$ governed by the DARE, yielding a formal bridge
between generative trajectory modeling and classical optimal control.

\subsection{Main Limitations}

MBD relies entirely on the analytical model and therefore cannot capture
multimodal, demonstration-informed behaviors---producing feasible but
potentially unnatural trajectories (e.g., ignoring preferred grasps or
learned obstacle-avoidance heuristics).
It also assumes a sufficiently accurate dynamics model, which may fail
under model mismatch, unmodeled contacts, or sim-to-real transfer, and
its importance-sampling score estimator suffers high variance when
$\lambda$ is small or trajectory dimension is large.
Our work targets these limitations: the learned prior restores
distributional richness, the adaptive $\lambda(x,t)$ weighting remains
robust under model imperfection, and the data score anchors the
denoising process to reduce variance.
MBD also provides no formal collision-avoidance guarantees; we plan to
add a CBF-derived score term~\cite{safeflow}.

% -------------------------------------------------------
\bibliographystyle{IEEEtran}
\bibliography{references}

\end{document}